\section{Related Work and Positioning}

This paper sits at the intersection of uncertainty-aware recognition, structured sequence decoding, and low-resource handwriting or glyph recognition. Prior work on calibration and uncertainty quantification has shown that posterior probabilities are often miscalibrated and that set-valued prediction can provide safer confidence statements~\cite{guo2017calibration,angelopoulos2021conformal}. Those results are important, but they are usually evaluated at the local prediction level rather than at the level of grouped or downstream inference.

Document analysis and handwriting recognition have long relied on structured decoding to preserve ambiguity long enough for lexical or language constraints to act~\cite{plamondon2000handwriting,graves2006ctc,scheidl2018wordbeam}. Our paper is close to that tradition in spirit, but it studies a narrower question: not whether a stronger recognizer wins outright, but whether preserving alternatives measurably changes what higher-level inference can recover in low-resource settings.

The benchmark design also matters. Omniglot and Kuzushiji-49 provide low-resource and manuscript-adjacent symbol corpora with very different visual statistics~\cite{lake2015omniglot,clanuwat2018kuzushiji}, while the Historical Newspapers and ScaDS.AI corpora provide the real grouped token alignments needed to test grouped rescue on actual sequences~\cite{historical_newspapers_gt,scadsai_grouped_words}. Within that landscape, the paper is intentionally positioned as a bounded empirical study of uncertainty propagation. It separates fully real evidence from synthetic-from-real evidence, and it does not claim decipherment, semantic recovery, or a universal law of uncertainty propagation.
